\begin{exblock}{复习题 3}
    \item 设 $f$ 与 $g$ 在 $(-\infty,+\infty)$ 上可导, 它们在 $-\infty$ 与 $+\infty$ 上分别存在有限的极限. 又设当 $x\in\mathbb{R}$ 时, $g'(x)\neq 0$. 证明: $\exists \xi\in(-\infty,+\infty)$, s.t.
    \[
    \frac{f(+\infty)-f(-\infty)}{g(+\infty)-g(-\infty)}=\frac{f'(\xi)}{g'(\xi)}.
\]

    \item 设 $f$ 与 $g$ 可导, 且对一切 $x$ 都有
    \[
    \begin{vmatrix}
    f(x) & g(x)\\
    f'(x) & g'(x)
\end{vmatrix}\neq 0.
\]
    证明: 在 $f$ 的任何两个不同零点之间, 至少有 $g$ 的一个零点.

    \item 设 $p$ 为一个实系数多项式, 再构造一个多项式
    \[
    q(x)=(1+x^{2})p(x)p'(x)+x(p(x)^{2}+p'(x)^{2}).
\]
    假设方程 $p(x)=0$ 有 $n$ 个大于 $1$ 的不同实根, 证明: 方程 $q(x)=0$ 至少有 $2n-1$ 个不同实根.

    \item 设 $n\in\mathbb{N}$, 且 $f(x)=\sum\limits_{k=1}^{n}c_{k}\mathrm{e}^{\lambda_{k}x}$, 其中 $\lambda_{1},\lambda_{2},\cdots,\lambda_{n}$ 为互不相等的实数, $c_{1},c_{2},\cdots,c_{n}$ 是不同时为 $0$ 的实数. 试问: 函数 $f$ 至多能有多少个实零点?

    \item 函数 $f:[a,b]\to\mathbb{R}$ 在 $[a,b]$ 上可导, 且 $f'(a)=f'(b)$. 证明: $\exists \xi\in(a,b)$, s.t.
    \[
    f'(\xi)=\frac{f(\xi)-f(a)}{\xi-a}.
\]

    \item 设函数 $f$ 在 $[a,+\infty)$ 上可导, $f(a)=0$, 且当 $x\geqslant a$ 时有 $|f'(x)|\leqslant |f(x)|$. 证明: $f=0$.

    \item 设 $f$ 在 $[0,+\infty)$ 上可导, 且 $0\leqslant f(x)\leqslant\frac{x}{1+x^{2}}$. 证明: $\exists \xi>0$, s.t.
    \[
    f'(\xi)=\frac{1-\xi^{2}}{(1+\xi^{2})^{2}}.
\]

    \item 设函数 $f$ 在 $[0,1]$ 上连续, 在 $(0,1)$ 内可导, 且 $f(1)-f(0)=1$. 证明: 对于 $k=0,1,\cdots,n-1$, $\exists \xi_{k}\in(0,1)$, s.t.
    \[
    f'(\xi_{k})=\frac{n!}{k!(n-1-k)!}\xi_{k}^{k}(1-\xi_{k})^{n-1-k}.
\]

    \item 设函数 $f$ 在 $[a,b]$ 上连续, 在 $(a,b)$ 内 $n$ 次可导. 又设 $a=x_{0}<x_{1}<\cdots<x_{n}=b$. 证明: $\exists \xi\in(a,b)$, s.t.
    \[
    \begin{vmatrix}
    1 & 1 & \cdots & 1\\
    x_{0} & x_{1} & \cdots & x_{n}\\
    \vdots & \vdots & & \vdots\\
    x_{0}^{n-1} & x_{1}^{n-1} & \cdots & x_{n}^{n-1}\\
    f(x_{0}) & f(x_{1}) & \cdots & f(x_{n})
\end{vmatrix}=\frac{1}{n!}f^{(n)}(\xi)\prod_{i>j}(x_{i}-x_{j}).
\]

    \item 设 $y_{1}=c>0$, $\frac{y_{n+1}}{n+1}=\ln\left(1+\frac{y_{n}}{n}\right)$, $n\in\mathbb{N}$. 求极限 $\lim\limits_{n\to+\infty}y_{n}$.

    \item 设函数 $f$ 在点 $x_{0}$ 有 $n$ 阶导数 ($n\in\mathbb{N}$). 证明:
    \[
    f^{(n)}(x_{0})=\lim_{h\to 0}\frac{1}{h^{n}}\sum_{k=0}^{n}(-1)^{n-k}C_{n}^{k}f(x_{0}+kh).
\]

    \item 求最大的数 $\alpha$ 及最小的数 $\beta$, 使满足下述不等式:
    \[
    \left(1+\frac{1}{n}\right)^{n+\alpha}\leqslant\mathrm{e}\leqslant\left(1+\frac{1}{n}\right)^{n+\beta},\quad n\in\mathbb{N}.
\]

    \item 设函数 $f$ 在点 $x_{0}$ 的开邻域内具有 $2$ 阶导数. 证明: 对充分小的 $h$, 存在 $\theta$, $0<\theta<1$, 使得
    \[
    \frac{f(x_{0}+h)+f(x_{0}-h)-2f(x_{0})}{h^{2}}=\frac{f''(x_{0}+\theta h)+f''(x_{0}-\theta h)}{2}.
\]

    \item 设函数 $f$ 在区间 $(a,+\infty)$ 上可导. 试用 L'Hospital 法则或 $\varepsilon\text{-}\Delta$ 法证明:
    \begin{enumerate}[label=(\arabic*)]
        \item 若 $\lim\limits_{x\to+\infty}f'(x)=0$, 则 $\lim\limits_{x\to+\infty}\frac{f(x)}{x}=0$;
        \item 若 $\lim\limits_{x\to+\infty}[f(x)+f'(x)]=r$, 则 $\lim\limits_{x\to+\infty}f(x)=r$.
    \end{enumerate}

    \item 设 $f$ 在 $(-r,r)$ 上有 $n$ 阶导数且 $\lim\limits_{x\to 0}f^{(n)}(x)=l$, $n\in\mathbb{N}$. 证明: $f^{(n)}(0)=l$.

    \item 设 $I$ 为开区间, 函数 $f$ 在 $I$ 上为凸函数 $\Longleftrightarrow \forall c\in I$, $\exists a\in\mathbb{R}$, s.t.
    \[
    f(x)\geqslant a(x-c)+f(c),\quad \forall x\in I.
\]
    对此作出几何解释.

    \item 设 $p$ 为多项式, 如果
    \[
    p'''(x)-p''(x)-p'(x)+p(x)\geqslant 0
\]
    在 $(-\infty,+\infty)$ 上成立. 证明: $p\geqslant 0$.

    \item 设 $f:[0,+\infty)\to\mathbb{R}$, 且 $\forall x\in[0,+\infty)$ 有 $x=f(x)\mathrm{e}^{f(x)}$. 证明:
    \begin{enumerate}[label=(\arabic*)]
        \item $f$ 是严格增的;
        \item $\lim\limits_{x\to+\infty}f(x)=+\infty$;
        \item $\lim\limits_{x\to+\infty}\frac{f(x)}{\ln x}=1$.
    \end{enumerate}

    \item 方阵 $\boldsymbol{A}=(a_{ij})$, $i,j=1,2,\cdots,n$, 其中一切元素均为正数, 其各行的和及各列的和均为 $1$. 设 $\boldsymbol{x}$ 是一个 $n$ 维列向量, 各分量均为正数. 令 $\boldsymbol{y}=\boldsymbol{A}\boldsymbol{x}$, 并设 $\boldsymbol{x}$ 与 $\boldsymbol{y}$ 的分量分别为 $x_{1},x_{2},\cdots,x_{n}$ 及 $y_{1},y_{2},\cdots,y_{n}$. 证明: $y_{1}y_{2}\cdots y_{n}\geqslant x_{1}x_{2}\cdots x_{n}$.

    \item 设 $a\geqslant 2$, $x>0$. 证明: $a^{x}+a^{\frac{1}{x}}\leqslant a^{\frac{x+1}{x}}$; 当且仅当 $a=2$ 与 $x=1$ 时, 式中等号成立.

    \item 微分方程
    \[
    \begin{cases}
    y''+p(x)y'+q(x)y=r(x),\quad x\in(a,b),\\
    y(a)=A,\\
    y(b)=B,
\end{cases}
\]
    其中 $q(x)<0$, $A,B$ 为常数, $p(x),q(x)$ 与 $r(x)$ 均连续. 如果这个方程在 $[a,b]$ 上有连续的解, 则解必是惟一的.

    \item 令 $p_{n}(x)=1+x+\frac{x^{2}}{2!}+\cdots+\frac{x^{n}}{n!}$, $n\in\mathbb{N}$. 证明:
    \begin{enumerate}[label=(\arabic*)]
        \item 当 $x<0$ 时, $p_{2n}(x)>\mathrm{e}^{x}>p_{2n+1}(x)$;
        \item 当 $x>0$ 时, $\mathrm{e}^{x}>p_{n}(x)\geqslant\left(1+\frac{x}{n}\right)^{n}$;
        \item 对一切实数 $x$, 有 $\mathrm{e}^{x}=\sum\limits_{n=0}^{\infty}\frac{x^{n}}{n!}$.
    \end{enumerate}

    \item 设 $f_{0}(x)=1$, $f_{k+1}(x)=xf_{k}(x)-f_{k}'(x)$. 证明:
    \begin{enumerate}[label=(\arabic*)]
        \item $f_{n}(x)$ 是首项系数为 $1$ 的 $n$ 次多项式;
        \item $f_{n}(x)$ 有 $n$ 个不同实根, 且关于 $0$ 对称地分布.
    \end{enumerate}

    \item 求函数 $f(x)=\mathrm{e}^{x}\left[\frac{1}{x}-\frac{\ln(x-1)}{x}\right]$ 在 $[2,4]$ 上的最大值.

    \item 设 $f(x)$ 满足 $f''(x)+f'(x)g(x)-f(x)=0$, 其中 $g(x)$ 为任一函数. 证明: 若 $f(x_{0})=f(x_{1})=0$ ($x_{0}<x_{1}$), 则 $f$ 在 $[x_{0},x_{1}]$ 上恒等于 $0$.

    \item 设 $f$ 在 $(-\infty,+\infty)$ 上 $2$ 阶可导, $f''>0$. 又设 $f(x_{0})<0$, 且 $\lim\limits_{x\to-\infty}f'(x)=\alpha<0$, $\lim\limits_{x\to+\infty}f'(x)=\beta>0$. 证明: 方程 $f(x)=0$ 在 $(-\infty,+\infty)$ 上恰有两个根.

    \item 设在有界闭区间 $[a,b]$ 上, 函数 $f$ 连续, $g$ 可导, $g(a)=0$, $\lambda\neq 0$ 为常数. 如果 $\forall x\in[a,b]$ 有
    \[
    |g(x)f(x)+\lambda g'(x)|\leqslant |g(x)|,
\]
    证明: $g(x)=0$, $\forall x\in[a,b]$.
\end{exblock}
